\documentclass[11pt,a4paper]{article}
\usepackage[utf8]{inputenc}
\usepackage[margin=1in]{geometry}
\usepackage{amsmath}
\usepackage{amssymb}
\usepackage{listings}
\usepackage{xcolor}
\usepackage{hyperref}
\usepackage{graphicx}
\usepackage{tikz}
\usepackage{float}

% Code listing style
\lstset{
    basicstyle=\ttfamily\small,
    keywordstyle=\color{blue},
    commentstyle=\color{green!60!black},
    stringstyle=\color{red},
    breaklines=true,
    frame=single,
    numbers=left,
    numberstyle=\tiny\color{gray},
    language=Python
}

\title{Comprehensive Notes: Cart-Pendulum with 2-DOF Planar Manipulator\\
\large Coupled Multi-Robot System in NVIDIA Isaac Sim}
\author{Script: \texttt{script\_cart\_pendulum\_manipulator\_basic\_run.py}}
\date{\today}

\begin{document}

% Cart–pendulum (small-angle, 2D horizontal) + impedance + muscle actuator
% Symbols:
%   x, x_ref, beta, F_imp, F_muscle, u \in R^2
%   m_c, m_b, l, g, tau_f scalars
%   K, D, M_ref \in R^{2x2} (typically diagonal)

\subsection*{Variables}
\begin{align*}
x &\in \mathbb{R}^2, & \dot{x} &\in \mathbb{R}^2,\\
\beta &\in \mathbb{R}^2, & \dot{\beta} &\in \mathbb{R}^2,\\
x_{ref} &\in \mathbb{R}^2, & \dot{x}_{ref} &\in \mathbb{R}^2,\\
F_{imp} &\in \mathbb{R}^2, & F_{muscle} &\in \mathbb{R}^2,\\
u &\in \mathbb{R}^2.
\end{align*}
Parameters:
\[
m_c,\; m_b,\; l,\; g,\; \tau_f \in \mathbb{R}, 
\qquad
K,\; D,\; M_{ref} \in \mathbb{R}^{2\times 2}.
\]

\subsection*{Impedance + muscle model}
\begin{align}
F_{imp} &= K\,(x-x_{ref}) + D\,(\dot{x}-\dot{x}_{ref}), \label{eq:impedance}\\
M_{ref}\,\ddot{x}_{ref} &= -F_{imp} + F_{muscle}, \label{eq:ref_dyn}\\
\tau_f\,\dot{F}_{muscle} &= u - F_{muscle}. \label{eq:muscle_filter}
\end{align}

\subsection*{Cart--pendulum dynamics (small-angle, spherical pendulum linearization)}
\begin{align}
m_c\,\ddot{x} &= F_{imp} + m_b\,g\,\beta, \label{eq:cart}\\
\ddot{\beta} &=
-\frac{(m_c+m_b)g}{l\,m_c}\,\beta
-\frac{1}{l\,m_c}\,F_{imp}. \label{eq:pend}
\end{align}

\subsection*{First-order state-space form}
Define the state:
\[
z =
\begin{bmatrix}
x\\ \dot{x}\\ \beta\\ \dot{\beta}\\ x_{ref}\\ \dot{x}_{ref}\\ F_{muscle}
\end{bmatrix}
\in\mathbb{R}^{14}.
\]
Then:
\begin{align}
\dot{x} &= \dot{x}, \label{eq:ss1}\\
\ddot{x} &= \frac{1}{m_c}F_{imp} + \frac{m_b g}{m_c}\beta, \label{eq:ss2}\\
\dot{\beta} &= \dot{\beta}, \label{eq:ss3}\\
\ddot{\beta} &=
-\frac{(m_c+m_b)g}{l\,m_c}\beta
-\frac{1}{l\,m_c}F_{imp}, \label{eq:ss4}\\
\dot{x}_{ref} &= \dot{x}_{ref}, \label{eq:ss5}\\
\ddot{x}_{ref} &= M_{ref}^{-1}\left(-F_{imp} + F_{muscle}\right), \label{eq:ss6}\\
\dot{F}_{muscle} &= -\frac{1}{\tau_f}F_{muscle} + \frac{1}{\tau_f}u, \label{eq:ss7}
\end{align}
with the algebraic force law \eqref{eq:impedance}:
\[
F_{imp} = K\,(x-x_{ref}) + D\,(\dot{x}-\dot{x}_{ref}).
\]

\end{document}